\chapter[\emj{💸}~Bellman-Ford (Camino más Corto con Pesos Negativos)]{\emj{💸}~Bellman-Ford (Camino más Corto con Pesos Negativos)}

\begin{dsaquees}

Buscar la ruta más barata entre ciudades 🏙️, pero donde algunas
carreteras te PAGAN por pasar (peso negativo), por ejemplo un descuento
o un reembolso.

Dijkstra se rompe con esos reembolsos porque asume que avanzar siempre
"cuesta". Bellman-Ford no asume nada: repite el cálculo muchas veces,
relajando todas las carreteras una y otra vez, hasta que ninguna ruta
puede mejorar más.

Además detecta las trampas: si existe un bucle que te paga infinito
(ciclo negativo), Bellman-Ford lo denuncia.

\end{dsaquees}

\begin{dsaotraforma}

Bellman-Ford calcula la distancia más corta desde un nodo origen a
todos los demás en un grafo dirigido y ponderado, admitiendo pesos
negativos (algo que Dijkstra no permite).

Su idea central es la "relajación" (relaxation): recorrer todas las
aristas $V-1$ veces, actualizando \verb|dist[v]| cada vez que
\verb|dist[u] + peso(u,v) < dist[v]|. Tras $V-1$ pasadas, las distancias
son óptimas si no hay ciclos negativos.

Una pasada extra ($V$-ésima): si alguna arista todavía se puede relajar,
existe un ciclo de peso negativo alcanzable desde el origen.

\end{dsaotraforma}

\begin{dsadiagrama}
\begin{verbatim}
Grafo dirigido con peso negativo:

   (4)      (-2)
 A ───> B ────> C
 │              ^
 │(5)           │(1)
 └──────> D ────┘

dist inicial: A=0, B=inf, C=inf, D=inf

Relajar aristas V-1 = 3 veces:
  A->B(4): B = 4
  B->C(-2): C = 2
  A->D(5): D = 5
  D->C(1): C = min(2, 6) = 2

Resultado: A=0, B=4, C=2, D=5

Pasada extra: ninguna arista mejora -> NO hay ciclo negativo.
\end{verbatim}
\end{dsadiagrama}

\begin{dsacomofunciona}
\begin{verbatim}
TABLA DE RELAJACIÓN (dist tras cada pasada), origen = A

Aristas: A→B(4)  B→C(-2)  A→D(5)  D→C(1)

           A     B     C     D
inicio  │  0    inf   inf   inf
pasada1 │  0     4     2     5      (A→B=4, B→C=2, A→D=5, D→C=min(2,6))
pasada2 │  0     4     2     5      (sin cambios → ya estable)
pasada3 │  0     4     2     5

Cada celda se actualiza solo si dist[u] + peso < dist[v].
Estable antes de V-1 pasadas → early exit posible.

DETECCIÓN de ciclo negativo (pasada extra):
  si ALGUNA arista aún cumple dist[u]+w < dist[v]  →  CICLO NEGATIVO
       A ──(1)──► B
       ▲          │
       └──(-3)────┘   suma del ciclo = 1 + (-3) = -2  →  negativo
\end{verbatim}
\end{dsacomofunciona}

\begin{lstlisting}
CÓMO FUNCIONA
1. dist[origen] = 0, el resto = infinito.
2. Repite V-1 veces: por CADA arista (u, v, peso), relaja:
   si dist[u] + peso < dist[v], entonces dist[v] = dist[u] + peso.
3. Pasada V-ésima: si alguna arista aún relaja -> ciclo negativo.

PSEUDOCÓDIGO
──────────────────────────────────────────────────────────────

función bellmanFord(V, aristas, origen):
    dist = arreglo lleno de INFINITO
    dist[origen] = 0
    REPETIR V-1 veces:
        PARA cada (u, v, w) en aristas:
            SI dist[u] != INF Y dist[u] + w < dist[v]:
                dist[v] = dist[u] + w
    // Deteccion de ciclo negativo
    PARA cada (u, v, w) en aristas:
        SI dist[u] != INF Y dist[u] + w < dist[v]:
            return "CICLO NEGATIVO"
    return dist

CÓDIGO C++
──────────────────────────────────────────────────────────────

#include <iostream>
#include <vector>
#include <climits>
using namespace std;

struct Edge { int u, v, w; };

bool bellmanFord(int V, vector<Edge>& edges, int src, vector<long long>& dist) {
    dist.assign(V, LLONG_MAX);
    dist[src] = 0;

    // V-1 pasadas de relajación
    for (int i = 0; i < V - 1; i++) {
        for (auto& e : edges) {
            if (dist[e.u] != LLONG_MAX && dist[e.u] + e.w < dist[e.v])
                dist[e.v] = dist[e.u] + e.w;
        }
    }
    // Pasada extra: si algo relaja, hay ciclo negativo
    for (auto& e : edges) {
        if (dist[e.u] != LLONG_MAX && dist[e.u] + e.w < dist[e.v])
            return false; // ciclo negativo detectado
    }
    return true;
}

COMPLEJIDAD (Big-O)
──────────────────────────────────────────────────────────────

  Caso        │ Tiempo      │ Espacio
  ────────────┼─────────────┼──────────────
  Tiempo      │ O(V * E)    │ O(V)

* Más lento que Dijkstra (O(E log V)) pero admite pesos negativos.
\end{lstlisting}

\begin{lstlisting}
✅ Optimización: si en una pasada completa NO se relajó ninguna arista,
   ya terminaste antes de tiempo (early exit).

💡 SPFA: una variante con Queue (Shortest Path Faster Algorithm) que
   suele ser más rápida en la práctica, aunque su peor caso sigue
   siendo O(V*E).

🔑 Aristas negativas SÍ, ciclos negativos NO: un peso negativo aislado
   está bien; lo que rompe todo es un CICLO cuya suma sea negativa.
\end{lstlisting}

\begin{dsaentrevistas}

✅ ¿Cuándo elegirlo sobre Dijkstra? Cuando haya pesos negativos o te
   pidan DETECTAR ciclos negativos. Es la respuesta esperada.

💡 "Cheapest Flights Within K Stops" (LeetCode 787): variante de
   Bellman-Ford limitando el número de pasadas a K+1.

⚠️ Inicializa con cuidado: nunca relajes desde un nodo con distancia
   infinita (evita overflow). Por eso el chequeo \verb|dist[u] != INF|.

🔑 Frase clave: "V-1 relajaciones porque el camino más corto tiene a lo
   sumo V-1 aristas". Dila y sumas puntos.

\end{dsaentrevistas}

\begin{dsaresumen}

• Camino más corto de una fuente a todos, CON pesos negativos.
• Relaja TODAS las aristas V-1 veces.
• Pasada extra detecta ciclos de peso negativo.
• O(V * E): más lento que Dijkstra pero más general.
• Usa Dijkstra si todos los pesos son positivos (más rápido).
• Clásico: rutas con reembolsos, arbitraje de divisas.
════════════════════════════════════════════════════════════════

\end{dsaresumen}
